% !TeX root = ../main.tex
\documentclass[../main.tex]{subfiles}

\begin{document}

\chapter{Espacio Vectorial Dual}


\section{Funcionales lineales y espacio dual}

Sea $\mathbb{V}(n,\mathbb{K})$ un espacio vectorial de dimensión $n$ sobre el campo $\mathbb{K}$. 
Sea $\{e_i\}_{i=1}^n$ una base de $\mathbb{V}(n,\mathbb{K})$, entonces si 
$A \in \mathbb{V}(n,\mathbb{K})$, en la base elegida se escribe como:
\begin{equation} \label{eq:vector_descomposicion}
    A = \sum_{i=1}^n A^i e_i,
\end{equation}
donde $A^i \in \mathbb{K}$ es la componente de $A$ en la dirección descrita por $e_i$.

Sea la \textbf{función lineal}
\begin{equation*}
    f : \mathbb{V} \longrightarrow \mathbb{K}
\end{equation*}

tal que
\begin{equation} \label{eq:linealidad}
    f(\alpha A_1 + \beta A_2) = \alpha f(A_1) + \beta f(A_2),
\end{equation}
$\forall \, \alpha, \beta \in \mathbb{K}, \;\; A_1, A_2 \in \mathbb{V}(n,\mathbb{K}).$

La función lineal nos dice que toma un vector $A \in \mathbb{V}(n,\mathbb{K})$ 
y lo lleva a un número en $\mathbb{K}$, esto es:
\begin{equation*}
    f(A) \in \mathbb{K}.
\end{equation*}

Usando \eqref{eq:vector_descomposicion} y la linealidad 
\eqref{eq:linealidad}, se encuentra:
\begin{align*}
    f(A) &= f\!\left(\sum_{i=1}^n A^i e_i \right) \\[6pt]
         &= f(A^1 e_1 + A^2 e_2 + \dots + A^n e_n) \\[6pt]
         &= A^1 f(e_1) + A^2 f(e_2) + \dots + A^n f(e_n) \\[6pt]
         &= \sum_{i=1}^n A^i f(e_i).
\end{align*}
\begin{equation*}
    \therefore \quad f(A) = \sum_{i=1}^n A^i f(e_i).
\end{equation*}
Por lo tanto, si conocemos los valores $f(e_i) \in \mathbb{K}$ 
para $i \in \{ 1,2,\dots,n\}$, sabremos el resultado de la operación de 
$f$ sobre cualquier vector $A \in \mathbb{V}(n,\mathbb{K})$.

El conjunto de las funciones lineales forma un \textbf{espacio vectorial}, 
puesto que una combinación lineal de dos funciones lineales es también una función lineal:
\begin{equation*}
    \big(\beta_1 f_1 + \beta_2 f_2\big)(A) 
    = \beta_1 f_1(A) + \beta_2 f_2(A).
\end{equation*}

Este espacio vectorial es llamado el \textbf{espacio vectorial dual} a $\mathbb{V}(n,\mathbb{K})$, 
y se denota por $\mathbb{V}^*(n,\mathbb{K})$ o simplemente $\mathbb{V}^*$.

Si la dimensión de $V$ es finita, entonces la dimensión de su dual es la misma:
\begin{equation*}
    \dim(\mathbb{V}) = \dim(\mathbb{V}^*).
\end{equation*}
    

\section{Base dual}

Puesto que $\mathbb{V}^*$ es un espacio vectorial, introducimos una base para $\mathbb{V}^*$ de la forma
\begin{equation}\label{eq:base_dual}
    \{E^1, E^2, \dots, E^n\} = \{E^i\}_{i=1}^n, 
    \qquad E^i \in \mathbb{V}^*.
\end{equation}

Como $E^i$ \textbf{es una función lineal}, queda completamente especificada si conocemos 
su acción sobre los elementos de la base $\{e_j\}$ de $\mathbb{V}$, es decir $E^i(e_j)$ para todo $i,j$.  

Elegimos la \textbf{base dual} de manera que satisfaga la condición:
\begin{equation}\label{eq:def_base_dual}
    E^i(e_j) = \delta^i{}_j.
\end{equation}

Cualquier función lineal $f$, llamado \textbf{vector dual}, 
se escribe en términos de la base dual \eqref{eq:base_dual} como
\begin{equation}\label{eq:vector_dual}
    f = \sum_{i=1}^n f_i E^i.
\end{equation}


\subsection{Acción del dual y producto interior}

La acción de $f$ sobre $A \in \mathbb{V}$ 
puede interpretarse como un \textbf{producto interno} 
entre las componentes $A^i$ y $f_i$. En efecto:

\begin{align*}
    f(A) &= f\!\left(\sum_{i=1}^n A^i e_i \right) \\[6pt]
    &= \sum_{i=1}^n A^i f(e_i) \\[6pt]
    &= \sum_{i=1}^n A^i \left( \sum_{j=1}^n f_j E^j(e_i) \right) \\[6pt]
    &= \sum_{i=1}^n \sum_{j=1}^n A^i f_j \, \delta^j{}_i \\[6pt]
    &= \sum_{i=1}^n A^i f_i. 
\end{align*}
\begin{equation}\label{eq:accion_de_f_sobre_A}
    f(A)= \sum_{i=1}^n A^i f_i. 
\end{equation}
Como se dijo, la expresión \eqref{eq:accion_de_f_sobre_A} corresponde a un 
\textbf{producto interior}, por lo que podemos escribir
\begin{align}
    \langle \; , \; \rangle &:\mathbb{V}^* \times \mathbb{V} \;\longrightarrow\; \mathbb{K}, \notag \\[6pt]
    &(f, A) \;\mapsto\;  f(A) = \sum_{i=1}^n A^i f_i. \label{eq:producto_interior}
\end{align}

y se denota como un producto interior:
\begin{equation*}
    \langle f, A\rangle := f(A).
\end{equation*}


\section{Índices covariantes y contravariantes}

En la expresión \eqref{eq:producto_interior} aparecen dos tipos de índices:  
uno arriba, llamado \textbf{contravariante}, y otro abajo, llamado \textbf{covariante}.  

La razón fundamental de esta notación ---un índice arriba y otro abajo--- 
es que la suma debe ser \textbf{invariante bajo transformaciones de coordenadas}.  
Es decir, la expresión \eqref{eq:producto_interior} debe ser la misma en todos los sistemas de coordenadas.

Para nuestro caso particular, el espacio vectorial es 
$\mathbb{V}(n,\mathbb{R}) = \mathbb{R}^3$. 
Los vectores base son la base coordenada $\{\vec{e}_i\}_{i=1}^3$ y, 
claramente, la acción de $f$ sobre $\vec{A} \in \mathbb{R}^3$ 
se interpreta, de acuerdo con \eqref{eq:producto_interior}, 
como el producto entre una \textbf{matriz fila $1 \times 3$} 
y una \textbf{matriz columna $3 \times 1$}.  

En efecto, las cantidades $A^i$ se pueden escribir como un vector columna:
\begin{equation*}
        A^i \;\longrightarrow\;
    \begin{pmatrix}
        A^1 \\
        A^2 \\
        A^3
    \end{pmatrix},
\end{equation*}

y las cantidades $f_i$ como un vector fila:
\begin{equation*}
    f_i \;\longrightarrow\;
    \begin{pmatrix}
        f_1 & f_2 & f_3
    \end{pmatrix}.    
\end{equation*}

Tal que
\begin{equation*}
    f(\vec{A}) =
    \begin{pmatrix}
        f_1 & f_2 & f_3
    \end{pmatrix}
    \begin{pmatrix}
        A^1 \\
        A^2 \\
        A^3
    \end{pmatrix}
    = \sum_{i=1}^3 A^i f_i.
\end{equation*}

Por lo tanto, conceptualmente, las cantidades con 
índices arriba y con índices abajo no son lo mismo, 
es decir:
\begin{equation*}
    T^i \;\neq\; T_i.
\end{equation*}

Así, se concluye que
\begin{equation*}
    \delta_{ij} \;\neq\; \delta^{ij} \;\neq\; \delta^j{}_i,
\end{equation*}

aunque los tres tengan el mismo valor en el sistema de coordenadas cartesianas.  
Sin embargo, cada uno representa un objeto distinto:

\begin{itemize}
    \item $\delta_{ij}$: es la \textbf{métrica euclideana} asociada a las coordenadas cartesianas  
    y no es invariante bajo transformaciones de coordenadas.

    \item $\delta^{ij}$: es la \textbf{métrica inversa} de la métrica euclideana  
    y tampoco es invariante bajo transformaciones de coordenadas.

    \item $\delta^j{}_i$: es la \textbf{identidad}, y sí es invariante bajo transformaciones de coordenadas.
\end{itemize}


\section{Base dual en $\mathbb{R}^3$}

Para $\mathbb{R}^3$, la base dual $\{E^i\}_{i=1}^3$ para \eqref{eq:producto_interior} satisface
\begin{equation}\label{eq:base_dual_R3}
    \delta^j{}_i = E^j(e_i) = \vec{E}^j \cdot \vec{e}_i.
\end{equation}

\begin{apuntebloque}{Definición}
Los \textbf{vectores base duales} para $\mathbb{R}^3$ se escriben de la forma
\begin{equation}\label{eq:suma_base_dual_R3}
    \vec{E}^i = \sum_{j=1}^3 g^{ij}\, \vec{e}_j,
\end{equation}
que satisface la condición \eqref{eq:base_dual_R3}. En efecto:
\begin{align*}
    \vec{E}^i \cdot \vec{e}_j 
        &= \sum_{k=1}^3 g^{ik}\, \vec{e}_k \cdot \vec{e}_j \\[6pt]
        &= \sum_{k=1}^3 g^{ik} g_{kj} \\[6pt]
        &= \delta^i{}_j.
\end{align*}
\end{apuntebloque}

Por lo tanto, sea $f \in \mathbb{V}^*(3,\mathbb{R})$ un vector dual de la forma
\begin{equation*}
    \vec{f} = \sum_{i=1}^3 f_i \, \vec{E}^i,
\end{equation*}
tal que, para $\vec{A} \in \mathbb{R}^3$, se tiene:
\begin{align*}
    f(A) 
        &= \vec{f} \cdot \vec{A} \\[6pt]
        &= \left( \sum_{i=1}^3 f_i \vec{E}^i \right) 
           \cdot \left( \sum_{j=1}^3 A^j \vec{e}_j \right) \\[6pt]
        &= \sum_{i=1}^3 \sum_{j=1}^3 f_i A^j 
           \big( \vec{E}^i \cdot \vec{e}_j \big) \\[6pt]
        &= \sum_{i=1}^3 \sum_{j=1}^3 f_i A^j \, \delta^i{}_j \\[6pt]
        &= \sum_{i=1}^3 f_i A^i.
\end{align*}

\begin{equation*}
    \therefore \quad f(A) = \sum_{i=1}^3 f_i A^i.
\end{equation*}



\subsection{Ejemplos de bases duales}

\begin{apuntebloque}{Ejemplos}
\textbf{a) Coordenadas Cartesianas:} Las bases coordenadas asociadas a las coordenadas cartesianas están dadas por 
$\hat{\imath}, \hat{\jmath}, \hat{k}$ y la métrica inversa es $\delta^{ij}$.  
Entonces, los vectores duales son:
\begin{align*}
    \vec{E}^1 &= \sum_{j=1}^3 \delta^{1j}\, \hat{e}_j 
    = \delta^{11}\hat{e}_1 + \delta^{12}\hat{e}_2 + \delta^{13}\hat{e}_3 
    = \hat{e}_1 = \hat{\imath}, \\[6pt]
    \vec{E}^2 &= \sum_{j=1}^3 \delta^{2j}\, \hat{e}_j 
    = \delta^{21}\hat{e}_1 + \delta^{22}\hat{e}_2 + \delta^{23}\hat{e}_3 
    = \hat{e}_2 = \hat{\jmath}, \\[6pt]
    \vec{E}^3 &= \sum_{j=1}^3 \delta^{3j}\, \hat{e}_j 
    = \delta^{31}\hat{e}_1 + \delta^{32}\hat{e}_2 + \delta^{33}\hat{e}_3 
    = \hat{e}_3 = \hat{k}.
\end{align*}

\textbf{Resumiendo:}
\begin{equation*}
\left.
\begin{aligned}
    \vec{E}^1 &= \hat{\imath}, \\
    \vec{E}^2 &= \hat{\jmath}, \\
    \vec{E}^3 &= \hat{k}
\end{aligned}
\;\right\}
\quad \Rightarrow \quad \vec{E}^i = \hat{e}_i
\end{equation*}

En conclusión: 
\begin{equation*}
    \vec{E}^i = \hat{e}_i.
\end{equation*}

\textbf{b) Coordenadas Esféricas:} Usando los vectores base de las coordenadas esféricas y la métrica inversa, tenemos:
\begin{align*}
    \vec{E}^r &= \sum_{j=1}^3 g^{1j}\, \vec{e}_j 
              = g^{11}\vec{e}_1 + g^{12}\vec{e}_2 + g^{13}\vec{e}_3 
              = \vec{e}_r, \\[6pt]
    \vec{E}^\phi &= \sum_{j=1}^3 g^{2j}\, \vec{e}_j 
              = g^{21}\vec{e}_1 + g^{22}\vec{e}_2 + g^{23}\vec{e}_3 
              = \frac{1}{r^2 \sin^2\theta}\,\vec{e}_\phi, \\[6pt]
    \vec{E}^\theta &= \sum_{j=1}^3 g^{3j}\, \vec{e}_j 
              = g^{31}\vec{e}_1 + g^{32}\vec{e}_2 + g^{33}\vec{e}_3 
              = \frac{1}{r^2}\,\vec{e}_\theta.
\end{align*}

\textbf{Resumiendo:}

\begin{equation*}
\left.
\begin{aligned}
    \vec{E}^r     &= \vec{e}_r = \hat{r}, \\[6pt]
    \vec{E}^\phi  &= \frac{\vec{e}_\phi}{r^2 \sin^2\theta} = \frac{\hat{\phi}}{r \sin\theta}, \\[6pt]
    \vec{E}^\theta &= \frac{\vec{e}_\theta}{r^2} = \frac{\hat{\theta}}{r}.
\end{aligned}
\;\right\}
\end{equation*}

\end{apuntebloque}


\section{Derivadas de la base dual}

\begin{apuntebloque}{Proposición}
Las \textbf{bases duales} $\{\vec{E}^i\}_{i=1}^3$ satisfacen
\begin{equation*}
    \frac{\partial \vec{E}^i}{\partial x^j}
    \;=\;
    - \sum_{k=1}^3 \Gamma^i{}_{kj}\, \vec{E}^k .
\end{equation*}

\textbf{Demostración.} 

Partiendo de
\begin{equation*}
    \vec{E}^i \cdot \vec{e}_l = \delta^i{}_l,
\end{equation*}
derivando con respecto a $\dfrac{\partial}{\partial x^j}$:
\begin{align*}
    \frac{\partial}{\partial x^j}\big(\vec{E}^i \cdot \vec{e}_l\big) 
    &= \frac{\partial}{\partial x^j}\delta^i{}_l \\
    \Rightarrow\quad\frac{\partial \vec{E}^i}{\partial x^j} \cdot \vec{e}_l 
       + \vec{E}^i \cdot \frac{\partial \vec{e}_l}{\partial x^j}
    &= 0 \\
    \frac{\partial \vec{E}^i}{\partial x^j} \cdot \vec{e}_l 
       + \vec{E}^i \cdot \sum_{n=1}^3\Gamma^n{}_{lj}\;\vec{e}_n
    &=  \\
    \frac{\partial \vec{E}^i}{\partial x^j} \cdot \vec{e}_l 
       +  \sum_{n=1}^3\Gamma^n{}_{lj}\;\vec{E}^i \cdot\vec{e}_n
    &= \\
    \frac{\partial \vec{E}^i}{\partial x^j} \cdot \vec{e}_l 
       +  \sum_{n=1}^3\Gamma^n{}_{lj}\;\delta^i_n
    &= \\
    \frac{\partial \vec{E}^i}{\partial x^j} \cdot \vec{e}_l 
       +  \Gamma^i{}_{lj}
    &= \\
    \frac{\partial \vec{E}^i}{\partial x^j} \cdot \vec{e}_l 
       +  \sum_{k=1}^3\Gamma^i{}_{kj}\, \delta^k_l
    &= \\
    \frac{\partial \vec{E}^i}{\partial x^j} \cdot \vec{e}_l 
       +  \sum_{k=1}^3\Gamma^i{}_{kj}\;\vec{E}^k \cdot\vec{e}_l
    &= \\
    \left( \frac{\partial \vec{E}^i}{\partial x^j} 
       +  \sum_{k=1}^3\Gamma^i{}_{kj}\;\vec{E}^k \right) \cdot\vec{e}_l
    &= \\
\end{align*}

Como los $\{\vec{e}_i\}_{i=1}^3$ son linealmente independientes, se obtiene:
\begin{equation*}
    \frac{\partial \vec{E}^i}{\partial x^j} 
       +  \sum_{k=1}^3\Gamma^i{}_{kj}\;\vec{E}^k=0\,.
\end{equation*}
Por lo tanto:
\begin{equation*}
\frac{\partial \vec{E}^i}{\partial x^j}
    = - \sum_{k=1}^3 \Gamma^i{}_{kj} \, \vec{E}^k \,.
\end{equation*}

\hfill $\square$
\end{apuntebloque}


\subsection{Ejemplos de descenso de índices}

\begin{apuntebloque}{Ejemplo}

\textbf{a) Coordenadas Cartesianas:}  
En coordenadas cartesianas se tiene que $\vec{E}^i = \hat{e}_i$ y $\Gamma^k{}_{ij} = 0$, entonces
\begin{equation*}
    \frac{\partial \vec{E}^i}{\partial x^j}
    = - \sum_{k=1}^3 \Gamma^i{}_{kj} \, \vec{E}^k
    = 0.
\end{equation*}

\textbf{b) Coordenadas Esféricas:}  
Podemos obtener $\dfrac{\partial \vec{E}^i}{\partial x^j}$ de dos formas:  
la primera es derivar los $\vec{E}^i$ con respecto de las coordenadas $(r,\phi,\theta)$, la segunda es usar la proposición anterior.

\medskip
\textit{Caso $i=j=1$:}
\begin{align*}
    \frac{\partial \vec{E}^1}{\partial x^1} 
        &= \frac{\partial \vec{E}^r}{\partial r}
        = \frac{\partial \hat{r}}{\partial r} = 0, \\[6pt]
    \frac{\partial \vec{E}^1}{\partial x^1} 
        &= -\sum_{k=1}^3 \Gamma^1{}_{k1}\,\vec{E}^k
        = -\Gamma^1{}_{11}\vec{E}^1 - \Gamma^1{}_{21}\vec{E}^2 - \Gamma^1{}_{31}\vec{E}^3 = 0.
\end{align*}

\textit{Caso $i=1,\,j=2$:}
\begin{align*}
    \frac{\partial \vec{E}^1}{\partial x^2}
        &= \frac{\partial \vec{E}^r}{\partial \phi}
        = \frac{\partial \hat{r}}{\partial \phi}
        = \frac{1}{r}\,\vec{e}_\phi
        = \frac{r \sin^2 \theta}{r^2 \sin^2 \theta}\,\vec{e}_\phi
        = \frac{1}{r}\,\vec{E}^\phi, \\[6pt]
    \frac{\partial \vec{E}^1}{\partial x^2}
        &= -\sum_{k=1}^3 \Gamma^1{}_{k2}\,\vec{E}^k
        = -\Gamma^1{}_{12}\vec{E}^1 - \Gamma^1{}_{22}\vec{E}^2 - \Gamma^1{}_{32}\vec{E}^3 \\[6pt]
        &= -\big(0\cdot\vec{E}^1 + (-r^2\sin^2\theta)\vec{E}^2 + 0\cdot\vec{E}^3\big) \\[6pt]
        &= r\sin^2\theta\,\vec{E}^\phi.
\end{align*}

Se puede seguir calculando los $\dfrac{\partial \vec{E}^i}{\partial x^j}$ por ambos métodos.
El primero se mostró para verificar la proposición, pero usando directamente la fórmula se obtiene:
\begin{align*}
    \frac{\partial \vec{E}^r}{\partial r} &= 0, &
    \frac{\partial \vec{E}^\phi}{\partial r} &= -\frac{1}{r}\,\vec{E}^\phi, &
    \frac{\partial \vec{E}^\theta}{\partial r} &= -\frac{1}{r}\,\vec{E}^\theta, \\[6pt]
    \frac{\partial \vec{E}^r}{\partial \phi} &= r\sin^2\theta\,\vec{E}^\phi, &
    \frac{\partial \vec{E}^\phi}{\partial \phi} &= -\frac{1}{r}\,\vec{E}^r - \cot\theta\,\vec{E}^\theta, &
    \frac{\partial \vec{E}^\theta}{\partial \phi} &= \sin\theta\cos\theta\,\vec{E}^\phi, \\[6pt]
    \frac{\partial \vec{E}^r}{\partial \theta} &= r\,\vec{E}^\theta, &
    \frac{\partial \vec{E}^\phi}{\partial \theta} &= -\cot\theta\,\vec{E}^\phi, &
    \frac{\partial \vec{E}^\theta}{\partial \theta} &= -\frac{1}{r}\,\vec{E}^r.
\end{align*}

\end{apuntebloque}


\section{Métrica y descenso de índices}

El producto escalar entre los vectores $\vec{A}$ y $\vec{B}$, como se estudió, 
está dado por
\begin{equation*}
    \vec{A}\cdot \vec{B} = \sum_{i=1}^3 \sum_{j=1}^3 g_{ij} A^i B^j.
\end{equation*}

Esta expresión se puede reescribir como:
\begin{align*}
    \vec{A}\cdot \vec{B} 
    &= \sum_{i=1}^3 \sum_{j=1}^3 g_{ij} A^i B^j \\[6pt]
    &= \sum_{j=1}^3 g_{1j} A^1 B^j 
      + \sum_{j=1}^3 g_{2j} A^2 B^j 
      + \sum_{j=1}^3 g_{3j} A^3 B^j \\[6pt]
    &= A^1 \sum_{j=1}^3 g_{1j} B^j 
      + A^2 \sum_{j=1}^3 g_{2j} B^j
      + A^3 \sum_{j=1}^3 g_{3j} B^j.
\end{align*}

\noindent
Por lo tanto:
\begin{equation*}
    \vec{A}\cdot \vec{B} = 
    \begin{pmatrix}
        \sum_{j=1}^3 g_{1j} B^j &
        \sum_{j=1}^3 g_{2j} B^j &
        \sum_{j=1}^3 g_{3j} B^j
    \end{pmatrix}
    \begin{pmatrix}
        A^1 \\[4pt] A^2 \\[4pt] A^3
    \end{pmatrix},
\end{equation*}
la cual es un producto matricial entre una matriz de $1\times 3$ y una de $3\times 1$, 
análoga a la expresión
\begin{equation*}
    f(A) = \sum_{i=1}^3 f_i A^i =
    \begin{pmatrix}
        f_1 & f_2 & f_3
    \end{pmatrix}
    \begin{pmatrix}
        A^1 \\[4pt] A^2 \\[4pt] A^3
    \end{pmatrix}.
\end{equation*}

De acuerdo a lo anterior, la matriz $A^i$ se interpreta como un vector columna de orden $3\times 1$, entonces la expresión
\begin{equation*}
    \sum_{j=1}^3 g_{ij} B^j
\end{equation*}
debe ser un vector fila de orden $1\times 3$, al igual que la cantidad $f_i$.

\medskip
Es decir:
\begin{equation*}
    \sum_{j=1}^3 g_{ij} B^j =
    \begin{pmatrix}
        \sum_{j=1}^3 g_{1j} B^j &
        \sum_{j=1}^3 g_{2j} B^j &
        \sum_{j=1}^3 g_{3j} B^j
    \end{pmatrix}.
\end{equation*}

Si bien el índice $j$ es un índice mudo, en la expresión anterior el índice $i$ es libre (no está sumado) y aparece en posición inferior. 
De este modo, la cantidad
\begin{equation*}
    \sum_{j=1}^3 g_{ij} B^j
\end{equation*}
posee un índice abajo.

\begin{apuntebloque}{Definición}
Sea el tensor contravariante $T^i$, se define el tensor covariante $T_i$, 
asociado a $T^i$, como
\begin{equation*}
    T_i := \sum_{j=1}^3 g_{ij} T^j \,.
\end{equation*}
\end{apuntebloque}

Por lo tanto, el producto escalar entre los vectores $\vec{A}$ y $\vec{B}$ 
se puede escribir como
\begin{align*}
    \vec{A} \cdot \vec{B} 
    &= \sum_{i=1}^3 \sum_{j=1}^3 g_{ij} A^i B^j \\[6pt]
    &= \sum_{i=1}^3 A^i \left( \sum_{j=1}^3 g_{ij} B^j \right) \\[6pt]
    &= \sum_{i=1}^3 A^i B_i \\[6pt]
    &= B(A).
\end{align*}

También se puede escribir en la forma
\begin{align*}
    \vec{A} \cdot \vec{B} 
    &= \sum_{i=1}^3 \sum_{j=1}^3 g_{ij} A^i B^j \\[6pt]
    &= \sum_{j=1}^3 B^j \left( \sum_{i=1}^3 g_{ij} A^i \right) \\[6pt]
    &= \sum_{j=1}^3 B^j A_j \\[6pt]
    &= A(B).
\end{align*}

La métrica $g_{ij}$ permite definir un mapeo biyectivo entre 
$\mathbb{V}$ y $\mathbb{V}^*$ como
\begin{equation*}
    g : \mathbb{V} \;\longrightarrow\; \mathbb{V}^* \,,
\end{equation*}
donde $g_{ij}$ actúa como
\begin{equation*}
    T^i \;\longmapsto\; \sum_{j=1}^3 g_{ij} T^j \,.
\end{equation*}

Como el mapeo es biyectivo, entonces existe la inversa
\begin{align*}
    g^{-1} :&\; \mathbb{V}^* \;\longrightarrow\; \mathbb{V}, \\
    \qquad 
    &T_i \;\longmapsto\; \sum_{j=1}^3 g^{ij} T_j \,.
\end{align*}

La métrica permite pasar de vectores contravariantes a covariantes, 
y la métrica inversa realiza el proceso inverso.

\begin{apuntebloque}{Definición}
Sea el tensor covariante $T_i$, se define el tensor contravariante $T^i$, 
asociado a $T_i$, como
\begin{equation*}
    T^i := \sum_{j=1}^3 g^{ij} T_j \,.
\end{equation*}
\end{apuntebloque}

La norma de un vector $\vec{A}$ se puede reescribir como
\begin{equation*}
    \|\vec{A}\| 
    = \sqrt{\sum_{i=1}^3 \sum_{j=1}^3 g_{ij} A^i A^j} 
    = \sqrt{\sum_{i=1}^3 A^i A^i} \,.
\end{equation*}

\begin{apuntebloque}{Ejemplo}

\textbf{a) Coordenadas Cartesianas:}  
Sea el vector $\vec{A}$, con componentes $A^i$ y métrica $\delta_{ij}$, entonces:

\begin{equation*}
    A_i = \sum_{j=1}^3 g_{ij} A^j = \sum_{j=1}^3 \delta_{ij} A^j = A^i \,.
\end{equation*}

Resumiendo:
\begin{align*}
    A_x &= A^x, \\
    A_y &= A^y, \\
    A_z &= A^z.
\end{align*}

\end{apuntebloque}
\begin{apuntebloque}{}

\textbf{b) Coordenadas Esféricas:}  
Sea $\vec{A}$ un vector con componentes $(A^r, A^\phi, A^\theta)$, entonces $A_i$ toma la forma:

\begin{align*}
    A_r &= \sum_{j=1}^3 g_{1j} A^j = g_{11}A^r + g_{12}A^\phi + g_{13}A^\theta = A^r, \\[4pt]
    A_\phi &= \sum_{j=1}^3 g_{2j} A^j = g_{21}A^r + g_{22}A^\phi + g_{23}A^\theta = r^2 \sin^2\theta \, A^\phi, \\[4pt]
    A_\theta &= \sum_{j=1}^3 g_{3j} A^j = g_{31}A^r + g_{32}A^\phi + g_{33}A^\theta = r^2 A^\theta.
\end{align*}

Resumiendo:
\begin{align*}
    A_r &= A^r, \\
    A_\phi &= r^2 \sin^2\theta \, A^\phi, \\
    A_\theta &= r^2 A^\theta.
\end{align*}

\end{apuntebloque}


\section{Covarianza del producto escalar}

Hasta aquí se ha introducido un índice arriba y otro abajo.  
La idea principal de esto es que el producto entre vectores sea 
\textbf{covariante}, es decir, que sea el mismo en todos los sistemas de coordenados.  

Lo anterior significa que, si tenemos los sistemas coordenados 
$Ox^1x^2x^3$ y $O'x'^1x'^2x'^3$, entonces:

\begin{equation*}
    \sum_{i=1}^3 A^i B_i \;\Big|_{O}
    \;=\;
    \sum_{i=1}^3 A'^i B'_i \;\Big|_{O'} \,.
\end{equation*}

Esto será demostrado más adelante, cuando se estudien las \textbf{transformaciones de coordenadas}.

Se debe tener en cuenta que las componentes del vector $\vec{A}$ 
son distintas para sistemas de coordenadas diferentes.  

Por ejemplo:
\begin{equation*}
    A_x \neq A_r, 
    \qquad A_y \neq A_\phi, 
    \qquad A_z \neq A_\theta \,,
\end{equation*}
sin embargo, todas ellas están relacionadas a través de una 
\textbf{transformación de coordenadas}.


\newpage

\end{document}
